% Generated draft: requires mathematical and source review.
\documentclass[11pt]{article}
\usepackage[margin=25mm]{geometry}
\usepackage{fontspec}
\setmainfont{Latin Modern Roman}
\usepackage{amsmath,amssymb,mathtools,unicode-math}
\setmathfont{Latin Modern Math}
\usepackage{enumitem,needspace,xurl,hyperref}
\setlength{\parindent}{0pt}
\setlength{\parskip}{5pt}
\setlength{\emergencystretch}{3em}
\newcommand{\N}{\mathbb{N}}
\newcommand{\Z}{\mathbb{Z}}
\newcommand{\Q}{\mathbb{Q}}
\newcommand{\R}{\mathbb{R}}
\newcommand{\C}{\mathbb{C}}
\newcommand{\F}{\mathbb{F}}
\newcommand{\A}{\mathbb{A}}
\newcommand{\PP}{\mathbb{P}}
\newcommand{\E}{\mathbb{E}}
\newcommand{\eps}{\varepsilon}
\newcommand{\dd}{\,\mathrm{d}}
\begin{document}
\section*{UnsolvedMath ID: 30001665}
\textbf{Source problem code:} catalogue entry
\subsection{Definitions and mathematical statement}
For a positive integer \(n\), a tournament \(T=(V,A)\) on \(n\) vertices is a directed graph with \(|V|=n\), no loops, and exactly one of \((u,v)\) and \((v,u)\) in \(A\) for every two distinct vertices \(u,v\in V\). An orientation \(P\) of the path on \(n\) vertices has vertices \(p_1,\ldots,p_n\) and, for each \(1\leq i<n\), exactly one of the arcs \((p_i,p_{i+1})\) and \((p_{i+1},p_i)\), with no other arcs.

Interpret \(T\) containing \(P\) as ordinary directed subgraph containment: there is a bijection \(f:V(P)\to V(T)\) such that
\[
(u,v)\in A(P)\quad\Longrightarrow\quad(f(u),f(v))\in A(T).
\]
Thus the copy is spanning, and additional tournament arcs between its vertices are allowed. This containment convention is explicit here but is not specified in the supplied statement.

The literal unrestricted assertion attributed in the record to Rosenfeld is: for every positive integer \(n\), every tournament on \(n\) vertices contains every orientation of the path on \(n\) vertices in the above sense.

This is a transcription of the supplied target, not a verified formulation of Rosenfeld's conjecture. The record supplies no lower bound on \(n\), exceptional tournaments, or exceptional path orientations; whether any such qualifications are missing remains unresolved. Its literature assessment concerns Sumner's conjecture, so that assessment cannot establish the status or intended qualifications of the displayed target.
\subsection{Short English statement}
The supplied assertion asks whether every tournament contains, as a spanning directed subgraph, every orientation of a path with the same number of vertices. The intended range and any exceptions require review. The supplied dated literature assessment discusses Sumner's conjecture rather than the Rosenfeld statement recorded here.
\subsection{Sources}
\begin{itemize}
\item Supplied local UnsolvedMath registry record, permanent ID \(30001665\), source problem code \texttt{catalogue entry}, title \emph{Rosenfeld Conjecture for Oriented Paths}. The record attributes the conjecture by name to Rosenfeld, lists proposed year \(2011\), and supplies no separate proposer. Its status label is \texttt{partially\_solved}; this is supplied metadata, not an independently verified mathematical conclusion.
\item Supplied report citation: \emph{Combinatorics (2011)}, Oberwolfach Reports, pp.\ 5--83, DOI \texttt{10.4171/owr/2011/01}. \url{https://doi.org/10.4171/owr/2011/01}. The background describes Oberwolfach Reports as documenting mathematical workshops and their open-problem discussions. No precise page or passage for this statement is supplied.
\item Supplied literature source, labeled K\"uhn--Mycroft--Osthus: \url{https://doi.org/10.1112/plms/pdq035}.
\item Supplied literature source, labeled ``Recent finite-bound work'': \url{https://arxiv.org/abs/2608.11667}.
\item Literature assessment dated exactly \texttt{2026-08-21}: the supplied curation states that Sumner's conjecture is proved for sufficiently large \(n\), that its all-\(n\) statement remains open, and that current work still treats finite bounds. This is supplied dated curation derived from the OWR record, not independently verified. The source itself requests independent verification of time-sensitive status and cited claims. Its relevance to the Rosenfeld statement is unresolved.
\end{itemize}
\end{document}
